\section{Discussion}
\label{sec:discussion}

\subsection{Interpretation of the proposed mechanism}
\label{sec:mechanism}

Across the \sepa{} $\lambda$ grid, composite fidelity changes only modestly while SI and HM change substantially. In the factorial experiment, enabling separability increases Vigitel Macro-F1 by approximately 0.10, while anchor pressure contributes essentially nothing. We read the strongest empirically supported claim from this evidence narrowly: not that \sepa{} is a universally superior generator, but that, in one dataset, a transparent post-generation objective controls a utility-relevant dimension that the composite fidelity score does not capture. The factorial result strengthens this reading precisely because the Vigitel change tracks the separability factor and not the anchor factor -- a distinction that a simple before/after comparison of a single tuned configuration could not have established.

The COVID-19 results constrain this claim considerably. The fixed-$\lambda$ study suggests a possible average improvement, but its confidence interval is wide; in the predeclared factorial analysis, treatment effects are small relative to fold variation and none is significant. A scalar separability margin may simply be ineffective when minority mortality cases are clinically heterogeneous, or it may select easy but unrepresentative prototypes without our being able to detect this from aggregate Macro-F1 alone. Stronger separation can mean clearer signal, but it can equally mean the removal of difficult patients, comorbidity interactions, or genuine label ambiguity -- and the two are observationally similar from the outside. We take the comparative pattern across datasets, rather than either result in isolation, to be the more informative finding: separability control is not a fixed property of the method but interacts with how learnable and homogeneous the minority class already is.

\subsection{Alternative explanation: an artificially easier problem}
\label{sec:alternative}

The most important open question is whether \sepa{} recovers class-conditional signal or merely constructs an easier classification problem. Higher \tstr{} performance may arise because selection concentrates class-conditional signal that CTGAN generated but did not otherwise emphasise. It may instead arise because selection removes borderline observations, narrows minority modes, or amplifies a relationship peculiar to the training fold. Held-out-real testing rules out a purely synthetic self-consistency artefact, but it does not distinguish between these two mechanisms, because both would produce the same aggregate improvement in Macro-F1. Boundary-stratified error analysis, calibration, real-minority cluster retention, label-permutation controls, and explicit shortcut-detection diagnostics are needed to separate them; we treat this as the central unresolved question for the method (Section~\ref{sec:limitations}) rather than as a peripheral caveat.

\subsection{Minority preservation and clinical realism}
\label{sec:minority}

Preserving class prevalence is necessary but not sufficient. A synthetic dataset can reproduce 83 deaths while collapsing several distinct death phenotypes into a single mode. SI and HM detect geometric cohesion but cannot determine whether clinically important subgroups survive selection. Evaluation therefore needs to include class-conditional distributions and dependencies, minority-class precision and recall, AUPRC, subgroup coverage, calibration, and within-positive-class diversity -- none of which the present experiments report, precisely because SI, HM, and the composite fidelity score are all that the current pipeline exports.

The \dacpf{} results reinforce this reading. Disease-anchor filtering improves both available classifiers relative to standard CTGAN, whereas age-only and combined filtering do not; anchors may therefore be redundant, misdirected, or too aggressively weighted for some clinical constructs even when they help for others. Domain knowledge should constrain implausible combinations without forcing every synthetic patient toward a class centroid, and blinded clinician review of individual cases and cohort-level relationships would help detect incoherent clinical combinations that geometric scores cannot.

\subsection{Does the causal component add value?}
\label{sec:causal-discussion}

The present experiments do not identify causal effects, and we do not claim otherwise. Anchor--target associations can reflect confounding, selection, measurement, or coding artefacts. A DAG learned from 334 COVID-19 records would also be unstable without informative priors and bootstrap analysis. More specifically, the \dacpf{} results do not by themselves demonstrate value from causal knowledge \emph{as such}: disease anchors help in some comparisons, age and combined filters do not, and no association-matched non-causal control condition was included to isolate whether the clinical motivation behind the anchors, rather than the anchors' statistical content, is doing the work. The defensible interpretation is narrower than ``causal filtering helps'': causal hypotheses can supply auditable constraints and diagnostics, but their incremental value over matched associational controls, and their sensitivity to graph uncertainty, remain to be isolated.

\subsection{Privacy implications}
\label{sec:privacy-discussion}

Nearest-neighbour distance checks and the absence of exact duplicates are sanity checks, not evidence of anonymity. \sepa{} explicitly prefers candidates close to the real manifold on the realism term, which may improve utility while simultaneously increasing disclosure risk for the same reason. Membership- and attribute-inference attacks, rarity-stratified nearest-record distributions, and differentially private variants of the pipeline are therefore required before any release claim, and we do not make one here. Privacy should be treated as an optimisation axis that competes with realism and separability, not as an incidental consequence of synthesis that can be assumed and checked only informally.

\subsection{Statistical limitations}
\label{sec:stat-limitations}

Five-fold cross-validation supplies five performance estimates, but the training sets overlap, and the within-fold factorial conditions additionally share a generator and candidate pool. Blocking by fold is the appropriate design choice for these paired contrasts, but the resulting ANOVA $p$-values remain approximate rather than exact, because the folds are not independent replications in the classical sense. Repeated nested cross-validation across independent partition and generator seeds is required before the Vigitel effect can be treated as a stable, confirmatory result rather than a strong signal within one experimental configuration; the Breusch--Pagan result for Vigitel further supports the use of robust or bootstrap intervals in that confirmatory work.

The factorial endpoint also averages CatBoost and XGBoost within each fold. This simplifies attribution but can mask a treatment-by-classifier interaction -- for instance, a classifier that is more sensitive to synthetic-record geometry could dominate the averaged effect. A confirmatory model should retain classifier identity as a repeated factor and report classifier-specific contrasts alongside the pooled estimate. For the same reason, the Vigitel $p$-value reported in Section~\ref{sec:factorial-results} should be read as conditional on this specific experimental design -- two classifiers, one generator, one dataset sample, recorded seeds -- and not as a probability that \sepa{} will succeed on a new dataset or deployment.

Taken together, the discussion above answers the paper's research questions only provisionally. Augmentation alone does not resolve weak class geometry, although it can yield modest, classifier-dependent improvements. Separability-aware selection is strongly effective for Vigitel but not demonstrated for COVID-19, and the factorial design attributes the Vigitel gain specifically to separability rather than to anchors, with no corresponding attribution available for COVID-19. Whether the observed gain reflects genuine recovery of class-conditional signal, and whether it survives scrutiny on clinical realism, minority coverage, and privacy, remains open. Section~\ref{sec:limitations} converts these open points into a concrete set of studies rather than leaving them as unresolved caveats.
